Crime imposes significant social costs, straining the criminal justice system and increasing both public and private security expenditures. Given these stakes, crime consistently ranks among the public's top concerns and features prominently in political debate \citep{Scartascini, gallup2023}. Police forces play a central role in controlling crime and maintaining public safety. While there is extensive evidence on the effects of police resources on crime \citep[see, for instance,][]{Chalfin_JEL,Levitt_1997,DiTella}, we know much less about how the deployment of those resources, and the strategy changes enabled by different territorial assignments, shape police effectiveness. Studies of public service provision show that decentralization and local specialization can improve performance by aligning operations with local needs \citep[see, for instance,][]{besley2003centralized}. This suggests that policing models that foster local knowledge and allow officers to adapt practices to local conditions may be more effective. Beat policing reflects these insights. Rather than rotating officers across large areas, it permanently assigns each officer to a smaller territorial unit, enabling the accumulation of local knowledge and adaptation to local crime patterns and community dynamics \citep{cordner1995community,skogan1997comminuty}. This model has become increasingly common across police departments worldwide \citep{NASEM2018,PCSP2014}. Yet despite its widespread adoption, rigorous causal evidence on its effects on public safety and security perceptions remains scarce.

This paper combines rich administrative and survey data from Chile to provide causal evidence that beat policing can significantly improve public security and police effectiveness along multiple dimensions. We exploit quasi-random variation from the staggered adoption of a large reform, \emph{Plan Cuadrante}, which introduced a beat-policing model in urban areas by dividing police operational areas---typically municipalities---into smaller zones---quadrants---and permanently assigning officers to them. This permanent assignment was designed to help officers accumulate local knowledge of territorial features, crime patterns, and communities, and to adapt their practices accordingly. Using a difference-in-differences strategy, we show that adopting municipalities experienced significant declines in victimization, reported crime, and perceived insecurity.

\emph{Plan Cuadrante} was implemented in municipalities with at least 70\% urban population and rolled out between 1998 and 2013 across the 150 municipalities that met this condition. Its central objective was to reduce victimization and insecurity perceptions. Rather than changing police hierarchy or operational doctrine, the reform reorganized frontline deployment by subdividing police operational areas---typically municipalities---into quadrants and permanently assigning officers to each of them.
On average, municipalities were divided into seven quadrants, with size and boundaries defined by local police authorities based on road networks, geographic features, and available resources.

To examine the effects of \emph{Plan Cuadrante}, we use household-level data from the National Urban Security Survey (ENUSC), which reports crime victimization, security perceptions, trust in police, and private security investment for roughly 25,000 urban households across 101 municipalities each year. We complement these data with administrative records on arrests and crimes reported to the police. To leverage staggered adoption, we employ a difference-in-differences strategy comparing early and late adopters using the procedure developed in \citet{Callaway}. Analyses based on survey data focus on 46 urban municipalities that adopted the program between 2005 and 2012 and one never-treated municipality; together, they represent over one-quarter of Chile's population. Analyses based on administrative data use a broader sample of 96 municipalities that adopted between 2003 and 2013 and 196 never-treated municipalities, representing over half of Chile's population.

We find that households in adopting municipalities become 10 percentage points (36\%) less likely to experience crime in the previous 12 months, with effects appearing within one year and growing over time. Consistent with this finding, we document a 10\% decline in crimes reported to the police. These estimates may understate the reform's true impact because not all crimes are reported and, by bringing police closer to communities, the program could have increased reporting rates. Nevertheless, the alignment between survey-based and administrative estimates provides reassuring evidence of the reform's effectiveness.

Beyond reducing crime, the reform also reduced insecurity perceptions and private investment in security. This is noteworthy because actual crime and crime perceptions can diverge, potentially leading to suboptimal security expenditures. For example, \citet{Ajzenman_Pato} show that rising immigration in Chile between 2000 and 2016 did not increase crime but did increase insecurity perceptions.

A potential identification concern is that crime reductions in adopting municipalities could reflect displacement to neighboring areas rather than genuine effectiveness gains. Hot-spot policing, which concentrates police resources in high-crime locations, has been shown to generate this type of spatial displacement \citep{Blattman_Tobon_JEEA}. In our setting, spillovers are less likely because municipalities are large geographic areas and beat policing reorganized deployment across the full municipal territory rather than concentrating resources in selected locations. Nevertheless, we directly assess spatial spillovers through three complementary exercises. We find no evidence that \emph{Plan Cuadrante} displaced crime to neighboring municipalities, diverted police resources from control to treated areas, or generated spillovers in victimization and insecurity outcomes in the ENUSC sample. These results support the interpretation that improvements in treated municipalities reflect genuine effectiveness gains rather than crime displacement.

To improve public safety and security perceptions, police work must strengthen prevention and crime resolution. Both dimensions can be affected by police resources, territorial deployment, and policing strategy. In \emph{Plan Cuadrante}, the core beat-policing components were the reorganization of deployment and the operational strategy adjustments enabled by that reorganization. At the same time, some municipalities experienced sizeable increases in police resources, which raises the question of whether resource expansion alone can explain the effects we estimate.

During our study period, public expenditure on security was rising, and most municipalities experienced increases in police resources. Police resources were not the central component of \emph{Plan Cuadrante}, but adopting the model required a minimum operational capacity, which in some municipalities implied increases in officers and patrol vehicles above those observed elsewhere. Thus, part of the effects we document could reflect resource increases. However, changes in police resources alone do not explain our findings. Municipalities that adopted \emph{Plan Cuadrante} but experienced resource increases similar to those in control municipalities still show large and statistically significant improvements in public safety outcomes. Moreover, for most outcomes, these effects are close in magnitude to those in municipalities where resources increased more substantially. In addition, variation in police resources unrelated to \emph{Plan Cuadrante} adoption does not seem to generate comparable gains. Taken together, these patterns indicate that while resource increases helped enable implementation and may explain part of the effects in some municipalities, the beat-policing components of the reform were central drivers of the overall gains.

Against this background, the core institutional change introduced by \emph{Plan Cuadrante} was on the deployment margin. The reform permanently assigned officers to smaller, fixed quadrants while leaving national operational guidelines unchanged across treated and control municipalities throughout our study period \citep{Manual_Carabineros2018,DIPRES2007_largo}. The Chilean police is a highly hierarchical institution with centrally defined doctrine and limited officer discretion, so adoption did not involve a top-down redesign of policing strategy. Instead, any strategic adjustments likely emerged from the new deployment itself. By remaining in the same territory, officers accumulated local knowledge of the area and its community and could adapt day-to-day practices---patrol routes and prevention activities---to local conditions. In this sense, \emph{Plan Cuadrante} changed police work primarily through deployment, with strategy evolving endogenously within a common institutional framework.

The evidence supports this interpretation. Following the adoption of \emph{Plan Cuadrante}, households report higher police presence in their neighborhoods, a pattern that is consistent with both resource expansion and deployment changes. However, our broader results suggest effects beyond patrolling intensity alone. While earlier work finds that police presence can reduce crime while officers are present but may be attenuated by displacement across places or time windows \citep{MASTROBUONI2019, Blanes_Mastrobuoni_2018}, we find declines in victimization and reported crime with no evidence of spillovers to neighboring municipalities. In line with this pattern, households report better police performance and greater trust in police, and administrative data show that arrests increase even as crime declines, including in municipalities with resource increases similar to those in control municipalities. Taken together, these findings suggest that beat policing improved public security through stronger deterrence and more effective prevention, response, and crime resolution under stable territorial assignment.

%--- Contributions to the literature
Our findings contribute to two related strands of the literature on police organization and public safety. First, they provide some of the first causal evidence on the effects of beat policing on crime and security perceptions. Beat policing has become one of the most widely adopted policing models globally \citep{NASEM2018,PCSP2014}, yet robust causal evidence on its effectiveness remains scarce. We show that this model can generate substantial reductions in crime victimization, reported crime, and insecurity perceptions.

These results also speak to recent work on other policing models. Hot-spot policing reallocates police resources toward high-crime locations and has been shown to reduce crime in targeted areas, though often with displacement to nearby places \citep{DiTella, Blattman_Tobon_JEEA, Draca_2011, Weisburd_restat}. Beat policing differs fundamentally. Rather than concentrating police presence in selected locations---where the same officers need not return repeatedly---it provides continuous, territory-wide coverage through permanent officer-to-area assignment. This distinction matters in our setting, where many quadrants were not crime hot spots and the reform generated municipality-wide declines in crime and insecurity perceptions, with no evidence of spillovers to neighboring municipalities. Community policing, another model that has received growing attention in the economics literature, emphasizes citizen participation in defining policing priorities and, in some cases, in implementation itself. While it has been shown to increase trust in police in developed countries, evidence from low- and middle-income settings points to more limited effects on crime and trust \citep{blair2021, Tobon}. Beat policing also differs from that approach. By permanently assigning officers to stable territorial units, it strengthens local knowledge and brings police closer to communities without requiring direct citizen participation in strategy or operations. Our findings show that this model can reduce crime and insecurity perceptions while increasing trust in police institutions in a middle-income country.

A second contribution is to the broader literature showing that specific features of police organization matter for police effectiveness and public security outcomes. Research has shown that faster police response times improve crime clearance \citep{Blanes_Kirchmaier}, that the effectiveness of police presence depends on how this is organized and deployed rather than simply on their intensity \citep{facchetti2025, MASTROBUONI2019, Blanes_Mastrobuoni_2018}, that shifts in police prioritization across crime types shape criminal behavior \citep{Adda_etal_2014}, and that better information systems and technology can improve police performance \citep{Mastrobuoni_2020}. Earlier work also found that areas covered by smaller, more localized police departments outperform, in terms of security, those covered by larger centralized ones, hypothesizing that this might be related, among other factors, to closer citizen-police engagement \citep{ostrom1973, ostrom1978}. We add to this literature by studying a policing model that combines several of these margins within a single organizational reform. By assigning officers permanently to smaller territorial units rather than rotating them across larger ones, beat policing is designed to increase police visibility, strengthen knowledge of local conditions, and improve the capacity of officers to adapt their practices to local crime patterns. In line with these mechanisms, we show that the adoption of \emph{Plan Cuadrante} increased perceived police presence in neighborhoods and raised police effectiveness, as arrests increased even while victimization and reported crime declined. These organizational improvements translated into meaningful gains in public safety and security perceptions.

A third contribution concerns the literature on police resources, which documents elasticities of crime with respect to officers and expenditure ranging from $-0.1$ to $-2$, with larger effects in developed settings and for violent crimes \citep{Chalfin_JEL, Levitt_1997, DiTella}. We complement this body of work by showing that the quantity of resources is not the only margin that matters. How those resources are deployed matters as well. \emph{Plan Cuadrante} improved public safety outcomes even in municipalities that experienced no increase in police resources relative to controls, suggesting that reorganizing deployment can raise public safety and police effectiveness without expanding the resource base.

Finally, we also contribute to the literature on police practices, institutional trust, and citizen cooperation. A persistent challenge is that security perceptions often diverge from actual crime trends \citep{Ajzenman_Pato}, limiting the scope for improving both simultaneously. Rising crime erodes institutional trust \citep{Blanco_2025}, while police misconduct undermines cooperation \citep{ang_etal_2025}.\footnote{Racial disparities in police behavior are also well documented; see, e.g., \citet{Goncalves_etal_2025}, who develop a method to correct for selection bias in estimates of racial disparities in the use of force.} Our results show that beat policing can improve both dimensions---security and security perceptions---by bringing officers into sustained engagement with local communities through permanent territorial assignment. We document not only increases in police trust and evaluations of police performance, but also substantial reductions in private security investment, indicating that households revised their threat assessments in response to the program. 

%Finally, our findings speak to a broader question about how organizational design shapes public-service effectiveness. A growing literature shows that the spatial organization of service provision matters. In policing, we show that a model organizing officers around fixed territorial units---allowing them to build local knowledge, coordinate with local actors, and adapt practices accordingly---can substantially improve both crime control and institutional legitimacy. This aligns with evidence from health and education showing that bringing providers closer to the populations they serve can improve outcomes \citep{besley2003centralized, Bardhan, Bloom_2015, Basurto}. In a rigorously identified setting, we show that this type of organizational proximity generates meaningful gains in police effectiveness and citizen security. JORGE: Podemos o bien quitar todo el parrafo como dice el reviewer 1, o no quitarlo y simplemente no vincularlo a la literatura de salud y educacion. Otra cosa, la frase en la linea 1 que habla sobre service organization no la he cambiado.

The rest of the paper is organized as follows. Section \ref{institutions} presents background information about Chilean police and the program we study. Section \ref{identification} introduces the data and empirical strategy. Section \ref{results} presents the main results and discusses potential drivers behind our findings. Finally, section \ref{conclusions} concludes.